% Section 2 of the report, standalone so it compiles on its own (same preamble as section1.tex).
% To move it into the report: everything between \begin{document} and \end{document}.
% Labels from Section 1 (sec:commanding, lst:recipe, sec:auto-mode, sec:listening, sec:sdk-stand,
% sec:bridge), the later sensing section (sec:sensing) and Results (sec:results) show as ?? here.
\documentclass[11pt,a4paper]{article}
\usepackage[utf8]{inputenc}\usepackage[T1]{fontenc}\usepackage{lmodern}
\usepackage[margin=2.2cm]{geometry}
\usepackage{tikz}\usetikzlibrary{positioning,arrows.meta,calc,fit,decorations.pathreplacing,backgrounds}
\usepackage{booktabs,array,amssymb,caption,listings,xcolor,url,amsmath}
\numberwithin{figure}{section}\numberwithin{table}{section}
\AtBeginDocument{\numberwithin{lstlisting}{section}}
\captionsetup{font=small,labelfont=bf}

\definecolor{motion}{HTML}{DCE8F5}
\definecolor{percep}{HTML}{E3F1E0}
\definecolor{outside}{HTML}{F3F3F3}
\definecolor{accent}{HTML}{C0504D}

\tikzset{
  comp/.style={draw, rounded corners=3pt, thick, inner sep=6pt, align=center},
  prog/.style={draw, rounded corners=2pt, fill=white, font=\small, align=center, minimum height=8mm, inner sep=4pt},
  dev/.style={draw, rounded corners=3pt, fill=outside, font=\small, align=center, minimum height=9mm, inner sep=5pt},
  link/.style={-{Stealth[length=2.2mm]}, thick},
  wire/.style={thick},
  lbl/.style={font=\footnotesize, align=center},
}

\lstset{language=Python, basicstyle=\ttfamily\small, columns=fullflexible, keepspaces=true,
  commentstyle=\color{gray}\itshape, keywordstyle=\bfseries, frame=lines, xleftmargin=1em,
  showstringspaces=false}

\begin{document}
\setcounter{section}{1}

\section{The Interface}
\label{sec:interface}

This section describes the program interface that was built on the findings of
Section~\ref{sec:commanding}. Subsection~\ref{sec:goals} states what the interface was designed to do, and
Subsection~\ref{sec:structure} shows how it is put together. Subsections~\ref{sec:state} to~\ref{sec:stopping}
go through its parts in the order a program uses them: connecting and reading the robot's state, holding
the controller, standing up, walking and stopping. Subsection~\ref{sec:errors} lists how each silent failure
is handled, and Subsection~\ref{sec:testing} describes the fake robot used for testing.
Subsection~\ref{sec:choices} discusses the main design choices and the alternatives that were not chosen,
Subsection~\ref{sec:limitations} lists what the interface does not protect against, and
Subsection~\ref{sec:interface-summary} sums up.

\subsection{Design Goals}
\label{sec:goals}

The recipe in Listing~\ref{lst:recipe} makes the robot walk, but it leaves the program to get every
detail right. The robot gives no warning when something goes wrong: a command sent in the wrong mode or at
the wrong moment is simply ignored, or does the opposite of what was intended. The interface was designed
around five goals that follow from this:

\begin{enumerate}
  \item \textbf{No silent failures.} Every known reason for the robot to ignore a command is either fixed
    automatically, or the call stops with an error message that says what is wrong and what to do.
  \item \textbf{Every motion ends with the robot stopped.} This holds however the motion ends: when it is
    done, when it is interrupted, and when the program crashes.
  \item \textbf{State as plain values.} The robot's posture, battery level and position can be read at any
    time, like ordinary variables.
  \item \textbf{Calls that wait and report.} A call such as \texttt{walk(2.0)} returns when the motion is
    over, and reports how far the robot actually moved and why it stopped.
  \item \textbf{No ROS~2 knowledge needed.} A program that only walks the robot should not have to deal
    with topics, message types or the ROS~2 machinery behind them.
\end{enumerate}

\subsection{Structure}
\label{sec:structure}

The interface is a Python package called \texttt{lite3}, which runs on the perception computer. A program
uses it through one class, \texttt{Lite3}. Creating a \texttt{Lite3} object connects to the robot, and its
methods (\texttt{stand}, \texttt{walk}, \texttt{turn} and so on) control it. Behind the single class, the
code is split into small files by topic, as shown in Figure~\ref{fig:structure}.

\begin{figure}[htbp]
\centering
\begin{tikzpicture}[node distance=5mm and 6mm]
  \node[dev, minimum width=60mm] (user) {A program: \texttt{bot.stand()}, \texttt{bot.walk(2.0)}};
  % the class and its parts
  \node[prog, below=12mm of user, xshift=-36mm] (post) {Posture\\ \footnotesize heartbeat, stand, sit,\\ \footnotesize auto mode, tricks};
  \node[prog, right=of post] (mot) {Motion\\ \footnotesize walk, turn, steer,\\ \footnotesize halt};
  \node[prog, right=of mot] (core) {Robot\\ \footnotesize connection, state,\\ \footnotesize stop, close};
  \node[prog, below=5mm of post] (udp) {Commands\\ \footnotesize the short UDP\\ \footnotesize commands};
  \node[prog, below=5mm of core] (st) {State\\ \footnotesize what the state\\ \footnotesize numbers mean};
  \begin{scope}[on background layer]
    \node[comp, fill=percep, fit=(post)(mot)(core)(udp)(st), inner sep=8pt] (cls) {};
  \end{scope}
  \node[font=\small\bfseries, anchor=south west] at (cls.north west) {\texttt{Lite3} (perception computer)};
  \draw[link] (user) -- (user |- cls.north);
  % robot side
  \node[prog, below=24mm of udp] (jy) {Walking controller\\ (\texttt{jy\_exe})};
  \node[prog, below=24mm of st, minimum width=26mm] (br) {ROS~2 bridge};
  \begin{scope}[on background layer]
    \node[comp, fill=motion, fit=(jy), inner sep=6pt] (mc) {};
  \end{scope}
  \node[font=\small\bfseries, anchor=north] at (mc.south) {Motion computer};
  \draw[link] (udp) -- node[lbl, left] {short commands,\\ heartbeat (UDP)} (mc.north -| udp);
  \coordinate (mid) at ($(mot.south)!0.55!(mot.south |- br.north)$);
  \draw[link] (mot.south) -- (mid) -| node[lbl, pos=0.25, below] {velocity (ROS~2)} ([xshift=-5mm]br.north);
  \draw[link] ([xshift=5mm]br.north) -- node[lbl, right] {state, odometry\\ (ROS~2)} ([xshift=5mm]br.north |- st.south);
  \draw[link] ([yshift=2mm]br.west) -- node[lbl, above] {velocity (UDP)} ([yshift=2mm]br.west -| jy.east);
  \draw[link] ([yshift=-2mm]br.west -| jy.east) -- node[lbl, below] {state (UDP)} ([yshift=-2mm]br.west);
\end{tikzpicture}
\caption{The parts of the interface and the two ways its commands reach the robot. Short commands (posture,
mode, heartbeat) go straight to the walking controller over UDP; velocity and state go through the vendor's
ROS~2 bridge, as in Listing~\ref{lst:recipe}.}
\label{fig:structure}
\end{figure}

Commands reach the robot in the two ways found in Section~\ref{sec:commanding}. Posture, mode and heartbeat
are short commands sent straight to the walking controller over UDP. Each is twelve bytes: a command code,
a value (used for joystick positions) and a zero. Velocity goes through the vendor's ROS~2 bridge, which also
brings the robot's state back. The codes the interface uses are listed in Table~\ref{tab:codes}. Each is
marked with where it came from, since some were recorded from the app with tcpdump, as described in
Subsection~\ref{sec:listening}, and some were taken from a third-party project on GitHub and not trusted
until they had been seen working on this robot. Besides these, the interface can play six built-in tricks,
also from the third-party project, of which only the dance has been run on this robot.

\begin{table}[htbp]
\centering
\caption{The short commands the interface sends to the walking controller, and where each was found.}
\label{tab:codes}
\begin{tabular}{@{}lll@{}}
\toprule
Code & Meaning & Found \\
\midrule
\texttt{0x21040001} & heartbeat, at least twice a second & recorded from the app \\
\texttt{0x21010202} & stand up or lie down (switches between the two) & recorded from the app \\
\texttt{0x21010C03} & auto mode: follow velocity commands & recorded from the app \\
\texttt{0x21010C02} & manual mode: follow the app's joysticks only & third-party \\
\texttt{0x21010D05}, \texttt{0D06} & enter and leave posture mode & recorded from the app \\
\texttt{0x21010130} & body pitch in posture mode (value: stick position) & recorded from the app \\
\bottomrule
\end{tabular}
\end{table}

Posture mode is a mode for tilting the body while standing still, used to point the depth camera up or
down. It mattered for the design because, in posture mode, the robot reports that it is standing but
ignores every velocity command.

\subsection{Connecting and Reading State}
\label{sec:state}

When a \texttt{Lite3} object is created, it first selects the messaging layer that ROS~2 runs on. ROS~2
can use several such layers, and with the default one the program found none of the robot's topics and
waited forever without an error. The interface therefore always selects the one the robot's own programs
use, before ROS~2 is loaded.

The object then starts listening to the topics the bridge publishes: the robot's state, its odometry (the
robot's own estimate of how far it has moved, calculated from the leg movements), the body's tilt, the two
ultrasonic distance sensors and the positions of the app's joysticks. A background thread, a part of the
program that runs alongside the rest, reads these topics 25 times a second and keeps only the newest
message from each. The state is then available as plain values such as \texttt{bot.battery},
\texttt{bot.pose} and \texttt{bot.standing}, which meets the third and fifth goals: a program reads them
like variables, without writing any ROS~2 code. If no state arrives within ten seconds, creating the object
fails with a message that names the likely cause, which is that the bridge is not running.

The most important value is a single state number that the walking controller reports. It says whether the
robot is standing, lying, or on its way between the two. Its values, found by watching it while the robot
was used, are listed in Table~\ref{tab:states}. The battery level is not part of the bridge's original
messages; the bridge was changed to pass it on, since the robot refuses to stand on a low battery without
saying why.

\begin{table}[htbp]
\centering
\caption{The values of the robot's state number that the interface relies on.}
\label{tab:states}
\begin{tabular}{@{}ll@{}}
\toprule
Value & Meaning \\
\midrule
6  & standing \\
1  & lying down, after a command to lie down \\
8  & lying down, no heartbeat: commands are ignored \\
98 & lying down, just after the robot was switched on \\
9  & on the way from 8 to 1; not safe to command \\
\bottomrule
\end{tabular}
\end{table}

\subsection{Holding the Controller}
\label{sec:heartbeat}

As found in Subsection~\ref{sec:listening}, the walking controller ignores network commands unless a
heartbeat keeps arriving. When a \texttt{Lite3} object is created, it starts a background thread that sends
the heartbeat twice a second, the same rate as the app, until the object is closed.

This has a consequence for safety. With the heartbeat running, the robot obeys the program with no app
connected, so the program, and the person running it, take the app's place as the stop. The app is not
made useless by this: touching one of its joysticks still stops the robot, as described in
Subsection~\ref{sec:stopping}.

An earlier version only sent the heartbeat when a program asked for it, and left it to the app otherwise.
This was changed after a test showed what the robot does when the heartbeat stops in the middle of a walk:
it keeps walking for a few seconds and then lies down where it is, which nearly tipped it over (the test is
described in Section~\ref{sec:results}). With the heartbeat left to the app, this would happen if the app
were switched off, or its battery ran out, while a program was driving the robot. With the program holding
its own heartbeat, the heartbeat only stops when the program itself stops, and then its velocity commands
stop with it.

\subsection{Posture}
\label{sec:posture}

Standing up looks like a single command, but Section~\ref{sec:commanding} left two problems with it. The
stand command switches between standing and lying, so sending it to a standing robot lays it down. And the
first stand command after the robot is switched on is ignored. Listing~\ref{lst:stand} shows how
\texttt{stand()} deals with both.

\begin{lstlisting}[float=htbp, captionpos=b, label=lst:stand,
  caption={How \texttt{stand()} works, simplified. The robot is only toggled from a settled lying state, and the toggle is sent once more if it was swallowed.}]
def stand():
    if state() == STANDING:
        return                                   # already up: sending would lie it down
    if state() not in LYING:
        raise Lite3Error('on the way between postures, not toggling')
    check_battery()                              # it refuses below 20 % without a word
    for attempt in (1, 2):
        send(STAND_TOGGLE)
        if wait_until_standing(give_up_if_still_lying_after=2.0):
            return
        # still lying: the first toggle after power-on only woke the controller
    raise Lite3Error('did not stand up - is the app or the heartbeat running?')
\end{lstlisting}

The method first reads the state number. If the robot is already standing, nothing is sent, so calling
\texttt{stand()} twice is harmless. If the robot is neither standing nor settled lying down, nothing is sent
either, since a toggle in the middle of a movement could lay it down. Only from a settled lying state is the
toggle sent. If the robot is still lying two seconds later, the first toggle was swallowed, and it is sent
once more. A shorter wait would risk sending the second toggle to a robot that was in fact getting up,
which would lay it down again. The method \texttt{sit()} works the same way in reverse: it stops any
movement, and toggles only a robot that is standing.

Auto mode is handled in a simpler way. As found in Subsection~\ref{sec:auto-mode}, the robot does not
report which mode it is in. The interface therefore never tries to find out. Every motion call leaves
posture mode and switches on auto mode before it moves, whether or not this is needed.

\subsection{Motion}
\label{sec:motion}

All motion goes through one loop, shown in Figure~\ref{fig:loop}. The calls \texttt{walk},
\texttt{strafe} (sideways), \texttt{turn} and \texttt{steer} differ only in what they do inside one cycle
of it. Before the loop starts, the robot must be standing, the battery high enough and the odometry
arriving; if not, the call stops with an error and the robot does not move. The loop then runs every
50~milliseconds. In each cycle it first runs the safety checks, which are the same for every call, and then
sends the call's velocity and checks whether the goal has been reached.

\begin{figure}[htbp]
\centering
\begin{tikzpicture}[node distance=6mm, box/.style={prog, minimum width=46mm}]
  \node[box] (in) {Checks before moving\\ \footnotesize standing, battery, odometry};
  \node[box, below=of in] (auto) {Leave posture mode,\\ switch on auto mode};
  \node[box, below=of auto] (safe) {Safety checks\\ \footnotesize app joysticks, cancel, odometry,\\ \footnotesize tilt, time limit};
  \node[box, below=of safe] (cyc) {One cycle of the call\\ \footnotesize send velocity, wait 50\,ms,\\ \footnotesize goal reached?};
  \node[box, below=10mm of cyc, fill=percep] (halt) {Stop: zero velocity\\ \footnotesize sent 20 times};
  \node[box, below=of halt] (out) {Return the report\\ \footnotesize reason, start, end, distance};
  \draw[link] (in) -- (auto);
  \draw[link] (auto) -- (safe);
  \draw[link] (safe) -- node[lbl, right] {all clear} (cyc);
  \draw[link] (cyc.west) -- ++(-8mm,0) |- node[lbl, pos=0.25, left] {not yet} (safe.west);
  \draw[link] (cyc) -- node[lbl, right] {goal reached} (halt);
  \draw[link] (safe.east) -- ++(10mm,0) |- node[lbl, pos=0.4, right] {reason\\ to stop} (halt.east);
  \draw[link] (halt) -- (out);
  \draw[link, dashed, accent] ([xshift=34mm,yshift=-6mm]cyc.east) node[lbl, above, text=black] {crash or\\ Ctrl-C} |- ([yshift=-2mm]halt.east);
\end{tikzpicture}
\caption{The one motion loop behind \texttt{walk}, \texttt{strafe}, \texttt{turn} and \texttt{steer}. The
loop runs every 50\,ms. However it ends, also by a crash, the robot is sent zero velocity before the call
returns or the error passes on.}
\label{fig:loop}
\end{figure}

The safety checks stop the motion if a person moves the app's joysticks, if another part of the program
cancels the motion, if the odometry has stopped arriving for half a second, if the body is tilted more than
30 degrees, or if the call has run for 30 seconds. Each of these gives a different reason, which is returned
with the result. A walk that stops early is not treated as an error: the program reads the reason, as in
Listing~\ref{lst:walk}.

\begin{lstlisting}[float=htbp, captionpos=b, label=lst:walk,
  caption={A program using the interface. Stopping short is reported, not raised; a refusal to start is raised.}]
with Lite3() as bot:                  # connects, starts the heartbeat, auto mode
    bot.stand()
    try:
        r = bot.walk(2.0, speed=0.3)
        print('moved %.2f m: %s' % (r['moved'], r['reason']))   # e.g. 'target reached'
        bot.turn_deg(90)              # positive is left
    except Lite3Error as e:
        print('refused:', e)          # the message says what to do
    finally:
        bot.sit()
\end{lstlisting}

The velocity a call sends is checked before it is sent. A walking speed above 0.6 metres per second is
refused, rather than quietly reduced. In \texttt{steer}, where an algorithm computes a new velocity every
cycle, the velocity is held within the limits, and a value that is not a valid number is sent as zero. This
last case matters more than it seems: limiting such a value with the usual minimum and maximum functions in
Python can give the highest allowed speed instead of a stop.

Turning has one more detail. The robot keeps turning for a short moment after it is told to stop, so
\texttt{turn()} sends the stop slightly before the target is reached, waits for the robot to settle, and
reports where it actually ended up. The angle is added up cycle by cycle, so turns of more than half a
revolution are counted correctly. Walking forward also stops for obstacles seen by the depth camera; how the
camera's view is turned into distances is described in Section~\ref{sec:sensing}.

\subsection{Stopping Safely}
\label{sec:stopping}

The second goal, that every motion ends with the robot stopped, is met in four layers. First, the motion
loop sends zero velocity on every way out, as shown in Figure~\ref{fig:loop}. In Python this is a
\texttt{finally} block, which runs even when the program crashes inside the loop. The zero velocity is sent
twenty times over 0.4 seconds, since the bridge's channel does not guarantee that a single message arrives.

Second, pressing Ctrl-C in the terminal runs an emergency stop: it cancels any running motion and sends
zero velocity. While the stop runs, further presses of Ctrl-C are ignored, so pressing it twice cannot cut
the stop short. Another part of the program, such as a web page used to control the robot, can stop a
motion at any time with \texttt{cancel()}.

Third, when the app is connected, it is a stop of its own. As soon as one of its joysticks is touched, the app sends
the command for manual mode, and from then on the robot ignores velocity from the program and follows the
joysticks instead. Touching a joystick therefore works as an emergency stop that needs no button. The
interface is built not to fight this. It switches on auto mode only when a motion call starts, never during
one, and the safety check described in Subsection~\ref{sec:motion} sees the joystick move and ends the call
with the reason that the app took over. A program that wants to move the robot again has to start a new call.

Fourth, closing the interface is done in a fixed order. The robot is stopped first, and only then is the
heartbeat stopped, since a robot that loses its heartbeat while it is still moving can tip over. The
ROS~2 parts are then shut down in an order that avoids crashes found during development. The object is meant
to be used in a \texttt{with} block, as in Listing~\ref{lst:walk}, which closes it in this order however the
block is left. Beyond these, the walking controller stops the robot by itself when velocity commands stop
arriving, so a program that is killed outright still leaves the robot standing still.

\subsection{Errors}
\label{sec:errors}

Table~\ref{tab:errors} lists the ways the robot fails silently that were found in
Section~\ref{sec:commanding} and while building the interface, and what the interface does about each. All
refusals are raised as one type of error, \texttt{Lite3Error}, with a message that names the cause.

\begin{table}[htbp]
\centering
\caption{Silent failures of the robot, and how the interface handles them.}
\label{tab:errors}
\begin{tabular}{@{}p{0.42\linewidth}p{0.52\linewidth}@{}}
\toprule
The robot & The interface \\
\midrule
ignores velocity in manual mode & switches on auto mode before every motion \\
ignores velocity in posture mode & leaves posture mode before every motion \\
ignores velocity when not standing & refuses to start, and says to call \texttt{stand()} \\
ignores everything without a heartbeat & sends its own heartbeat from the start \\
ignores the first stand after power-on & sends the stand command once more \\
lies down on a stand command when standing & sends nothing to a standing robot \\
refuses to stand on a low battery & refuses below 20\,\%, with the battery level \\
gives no state with the wrong ROS~2 setup & selects the right messaging layer; if no state comes, says which program to check \\
keeps walking on the last command after a crash & sends zero velocity on every way out \\
\bottomrule
\end{tabular}
\end{table}

\subsection{Testing Without the Robot}
\label{sec:testing}

The package includes a fake robot, which replaces the connection to the real one with an imitation whose
position follows the velocity it is sent, on a clock that only moves forward when the code waits. No message
leaves the computer, so a program built on the interface can be tested with no robot and no risk, and a
ten-second walk is checked in a fraction of a second. This is how the guarantee of Subsection~\ref{sec:stopping}
was checked, since making a program crash in the middle of a walk is not something to try on the real robot:
the tests check that every motion, however it ends, leaves zero velocity behind. The code described in the
later sections was tested against the same fake robot before it was run on the robot.

\subsection{Design Choices}
\label{sec:choices}

Three choices shaped the interface more than any other. Each had an alternative that was considered and not
chosen.

\paragraph{A Python library, not a ROS~2 node.} The obvious alternative was to extend the vendor's ROS~2
bridge, or to write a ROS~2 node beside it, offering stand, sit and walk as ROS~2 services that other
programs call over the network. This would have let several programs share the robot. It would also have
made every program that uses the robot a ROS~2 program, and split the robot's control across two processes:
if the program giving the orders crashed, the node carrying them out would keep running. As a library, the
interface runs inside the program that moves the robot. The stop guarantees of Subsection~\ref{sec:stopping}
then follow from this: when the program ends, in any way, its commands and its heartbeat end with it. The
cost is that only one program can control the robot at a time. Programs that do want ROS~2, such as the
navigation in the later sections, still have it: the interface gives them its ROS~2 node to add their own
topics to.

\paragraph{Calls that wait until the motion is done.} A call such as \texttt{walk(2.0)} returns only when the
motion is over. The alternative was calls that start a motion and return at once, leaving the program to
check on it later. Waiting calls keep programs simple, since they read in the order things happen, and they
give each motion a clear end, where the robot is stopped and the result reported. The cost is that a program
does only one motion at a time, and something else that should stop it, such as a button on a web page, has to
run in a separate thread and use \texttt{cancel()}. Navigation, described later, needs the opposite: it sends
a goal and keeps checking on it while the robot walks.

\paragraph{Fix what is safe to fix, refuse the rest.} For each silent failure, the interface either fixes it
or refuses with a message. It fixes a problem itself only when the fix is safe and is what the program
obviously meant: switching on auto mode, leaving posture mode, and sending a swallowed stand command again.
It refuses when a fix would hide something the person running the program should know, or would do something
they did not ask for. A walk on a robot that is lying down is refused rather than standing the robot up
first, since standing up is a large movement of its own. A speed above the limit is refused rather than
reduced, since a program that asks for too much has a bug that should be seen. A battery below 20~\% is
refused rather than tried, since the robot would refuse anyway, without saying why.

\subsection{Limitations}
\label{sec:limitations}

The interface does not protect against everything. The main limitations are:

\begin{itemize}
  \item \textbf{Only walking forward and turning look for obstacles.} A turn checks the side it turns to.
    Walking backward or sideways, and motion steered with \texttt{steer}, do not use the depth camera, and
    the camera only sees ahead. A program that uses
    these must check the way itself, as the navigation in the later sections does.
  \item \textbf{The stops are in software.} The only stops that do not depend on a program are touching the
    app's joysticks and the robot's power switch. With the app switched off, the program and the person
    running it are the only stop.
  \item \textbf{A command that does nothing is not noticed at once.} The robot does not confirm commands, and
    auto mode cannot be read back. If the robot ignores a motion, the call only notices that it has not moved
    when its time limit runs out.
  \item \textbf{Some checks have not happened on the robot.} The stop for a tilted body has never been
    triggered on the robot, and of the built-in tricks only the dance has been tried.
  \item \textbf{Losing the heartbeat while moving can tip the robot over.} Since the program holds the
    heartbeat, this only happens when the program itself stops, and then its velocity commands stop too.
    It is still the most serious way the robot can fail, and it is why the interface always stops the
    robot before it lets the heartbeat go.
  \item \textbf{One program at a time.} Two programs controlling the robot at once would send it conflicting
    commands, and nothing in the interface prevents it. The web panel must be stopped before another program
    is run.
\end{itemize}

\subsection{Summary}
\label{sec:interface-summary}

The interface turns the recipe of Listing~\ref{lst:recipe} into a single class that a program can use
without knowing how the robot is commanded underneath. Each silent failure found in
Section~\ref{sec:commanding} is either fixed automatically or reported with a message that says what to do,
and every motion ends with the robot stopped, also after a crash. Its main limits are that only walking
forward and turning look for obstacles, and that its stops are in software. The interface covers the robot as it came
from the manufacturer: posture, walking, turning, its state, the depth camera and the speaker. It is
published as a separate package, so that it can be used on any Lite3. Everything added to the robot later
in this project, such as the lidar and navigation, builds on this class rather than changing it.

\end{document}
